\section{Conclusiones}

\begin{itemize}
\item El paralelismo por cuerpo es natural para N-Body: las tareas son independientes dentro de un
paso y la única sincronización necesaria es la barrera entre fases y entre pasos.
\item La mejor configuración fue la versión C con \code{reduction}, \code{dynamic}, chunk 64 y 8
hilos: 0.069\,s, speedup de 5.17 y eficiencia de 0.65, con resultados numéricamente equivalentes a la
versión secuencial.
\item La estrategia contra la condición de carrera fue el factor más determinante:
\code{reduction} ($5.18\times$) frente a \code{atomic} ($0.65\times$) y \code{critical}
($0.14\times$).
\item El schedule y el chunk importan: \code{dynamic} con chunks de 8 a 64 fue el mejor en ambas
versiones, mientras que chunks de 512 y \code{guided} por defecto fueron los peores, especialmente con
la carga triangular de C.
\item Limitaciones: con $N = 5000$ y 10 pasos las corridas duran menos de 0.4\,s, por lo que existe
ruido de medición, y los núcleos heterogéneos del M2 Pro limitan la eficiencia con 8 hilos.
\end{itemize}
